\section*{Chapitre \ref{ch:b1:ocean-circulation} --- La circulation de l'océan}
% values: calc/ch29.py

\begin{solution}{exo:b1:ocean-circulation:1}
(a)~À \qty{4}{m/s} (première formule)~: $1000\,C_D = 0.29 + 0.775 + 0.481 = 1.546$,
$T = 1.3 \times \num{1.546e-3} \times 16 = \qty{0.0322}{Pa}$. À \qty{7}{m/s}
(seconde formule)~: $1000\,C_D = 0.60 + 0.497 = 1.097$, $T = \qty{0.0699}{Pa}$. À
\qty{15}{m/s}~: $1000\,C_D = 1.665$, $T = \qty{0.487}{Pa}$. (b)~$0.487/0.0699 =
6.97$, contre $(15/7)^2 = 4.59$~: le coefficient de traînée lui-même croît avec le vent,
parce qu'un vent plus fort soulève des vagues plus rugueuses qui agrippent mieux l'air.
\end{solution}

\begin{solution}{exo:b1:ocean-circulation:2}
(a)~Hémisphère Nord, \ang{90} à droite d'un vent soufflant vers l'est~: vers le
sud. (b)~Le vent souffle vers le sud-ouest~; à sa droite se trouve le nord-ouest.
(c)~Hémisphère Sud, \ang{90} à gauche d'un vent soufflant vers l'est~: vers le
nord. (d)~Le vent souffle vers le nord~; à sa gauche se trouve l'ouest, qui est
le large sur une côte ouest. Dans le cas (d) le transport pointe vers le large~: l'eau de surface
est remplacée par en dessous, et de l'eau froide remonte le long de la côte.
\end{solution}

\begin{solution}{exo:b1:ocean-circulation:3}
(a)~$31.7 + 3.5 - 17.9 = \qty{17.3}{\sverdrup}$, contre les
\qty{17.1}{\sverdrup} publiés~: les moyennes des trois pièces s'additionnent à la moyenne du retournement
aux arrondis et incertitudes près ($\pm\qty{3.2}{\sverdrup}$).
(b)~Au-dessus de \qty{1100}{m}~: $16.6 + 0.4 = \qty{17.0}{\sverdrup}$ vers le nord~; en dessous~:
$12.0 + 5.9 = \qty{17.9}{\sverdrup}$ vers le sud. Les deux s'équilibrent presque, comme ils le doivent
quand aucune eau ne s'accumule au nord de la section~; la petite différence est dans les
incertitudes. (c)~$\num{17.1e6} \times \num{3.156e7} = \qty{5.40e14}{m^3}$, soit
\qty{5.40e5}{km^3} par an~; $\num{1.3324e9}/(\num{5.40e5}) \approx 2470$ ans.
\end{solution}

\begin{solution}{exo:b1:ocean-circulation:4}
$D_E = 7.6 \times 10/\sqrt{\sin\varphi}$ et $V_0 = 0.0127 \times
10/\sqrt{\sin\varphi}$. À \ang{15}~: \qty{149}{m} et \qty{0.250}{m/s}
(\qty{2.5}{\percent} du vent)~; à \ang{45}~: \qty{90}{m} et \qty{0.151}{m/s}
(\qty{1.5}{\percent})~; à \ang{60}~: \qty{82}{m} et \qty{0.136}{m/s}
(\qty{1.4}{\percent}). Près de l'équateur $\sin\varphi \to 0$ et les deux formules
divergent~; plus fondamentalement $f \to 0$, si bien que la force de Coriolis nécessaire pour équilibrer
le frottement s'évanouit, et le temps $1/f$ dont l'équilibre a besoin pour s'installer devient plus long
que la durée pendant laquelle les vents restent stables. Ekman a restreint sa formule à $|\varphi| \geq
\ang{10}$.
\end{solution}

\begin{solution}{exo:b1:ocean-circulation:5}
(a)~$f = 2\Omega\sin\ang{11} = \qty{2.783e-5}{s^{-1}}$~; $M = 0.07/(\num{2.783e-5}) =
\qty{2515}{kg.m^{-1}.s^{-1}}$, \ang{90} à droite d'un vent soufflant vers l'ouest~:
vers le nord. (b)~$Q = (2515/1027) \times \num{5e6} = \qty{1.22e7}{m^3/s}$,
\qty{12.2}{\sverdrup}. (c)~L'estimation se situe dans l'incertitude de la mesure
directe ($12.0 \pm 5.5$) et près de la valeur issue des vents à long terme
($13.5 \pm 0.3$)~: le théorème rend compte du transport observé.
\end{solution}

\begin{solution}{exo:b1:ocean-circulation:6}
(a)~$f(\ang{2}) = \qty{5.09e-6}{s^{-1}}$~; $M = 0.05/(\num{5.09e-6}) =
\qty{9824}{kg.m^{-1}.s^{-1}}$, \qty{9.57}{m^2/s}~: à \ang{2}~N vers le nord
(à droite d'un vent vers l'ouest), à \ang{2}~S vers le sud (à gauche). (b)~Sur une bande
longue de \qty{1}{m} en \omterm{def:b1:spherical-astronomy:geographic}{longitude}, $2 \times 9.57 = \qty{19.1}{m^3/s}$ partent à travers les
deux \omterm{def:b1:spherical-astronomy:geographic}{latitudes}, distantes de \qty{444.8}{km}~: $w = 19.1/(\num{4.448e5}) = \qty{4.30e-5}{m/s}$,
environ \qty{3.7}{m} par jour. (c)~À l'équateur $f = 0$ et $|\vec M| = T/|f|$ est
infini~: le bilan d'Ekman permanent n'y existe pas, et les nombres ci-dessus, très
sensibles à la \omterm{def:b1:spherical-astronomy:geographic}{latitude} choisie, ne sont qu'un ordre de grandeur. Ce qui survit est
le signe~: de chaque côté le transport s'éloigne de l'équateur, si bien que l'eau de surface
diverge et que de l'eau froide monte le long de lui.
\end{solution}

\begin{solution}{exo:b1:ocean-circulation:7}
(a)~$f(\ang{15}) = \qty{3.775e-5}{s^{-1}}$~: $M = 0.06/(\num{3.775e-5}) =
\qty{1590}{kg.m^{-1}.s^{-1}}$ vers le nord~; $f(\ang{40}) =
\qty{9.375e-5}{s^{-1}}$~: $M = \qty{1280}{kg.m^{-1}.s^{-1}}$ vers le sud.
(b)~Entre les deux \omterm{def:b1:spherical-astronomy:geographic}{latitudes}, distantes de \qty{2780}{km}, $\partial M_y/\partial y =
-(1590 + 1280)/(\num{2.78e6}) = \qty{-1.03e-3}{kg.m^{-2}.s^{-1}}$~: d'après la
\cref{prop:b1:ocean-circulation:pumping}, $w_E = \qty{-1.005e-6}{m/s}$, une vitesse
descendante de \qty{31.7}{m} par an. (c)~$\num{1.005e-6} \times \num{2.78e6} \times
\num{5e6} = \qty{1.40e7}{m^3/s}$, \qty{14.0}{\sverdrup} pompés vers le bas.
\end{solution}

\begin{solution}{exo:b1:ocean-circulation:8}
(a)~$H = 1027 \times 3992 \times \num{17e6} \times 12 = \qty{8.36e14}{W}$,
\qty{0.84}{PW}. (b)~C'est \qty{70}{\percent} des \qty{1.2}{PW} mesurés~: la vraie
branche supérieure est plus chaude près de la surface que la moyenne de \qty{15}{\celsius}, et le
\omterm{def:b1:ocean-circulation:gyre}{gyre} porte aussi de la chaleur. (c)~$1027 \times 3992 \times \num{17e6} \times 3 =
\qty{2.1e14}{W}$, \qty{0.21}{PW}, un quart de (a). Seule l'eau qui devient froide
avant de revenir porte beaucoup de chaleur~: le transport de chaleur de l'Atlantique dépend de
la \omterm{def:b1:ocean-circulation:moc}{formation d'eau profonde}.
\end{solution}

\begin{solution}{exo:b1:ocean-circulation:9}
Soit $x$, $y$ et $1 - x - y$ les fractions d'AAIW, de NADW et d'AABW. Les bilans de sel et de chaleur,
mesurés à partir de l'AABW~:
\[
  -0.52\,x + 0.33\,y = 34.731 - 34.67 = 0.061,\qquad
  2.9\,x + 3.6\,y = 2.78 - 0.4 = 2.38 .
\]
Le déterminant est $-0.52 \times 3.6 - 0.33 \times 2.9 = -2.829$, et la règle de Cramer
donne $x = (0.061 \times 3.6 - 0.33 \times 2.38)/(-2.829) = 0.20$ et $y = (-0.52
\times 2.38 - 2.9 \times 0.061)/(-2.829) = 0.50$~: l'échantillon est à \qty{20}{\percent}
d'AAIW, \qty{50}{\percent} de NADW et \qty{30}{\percent} d'AABW, un point à l'intérieur du
triangle des trois eaux.
\end{solution}

\begin{solution}{exo:b1:ocean-circulation:10}
(a)~Avec $W = u + iv = V_0e^{i\pi/4}e^{(1+i)az}$, dont les parties réelle et imaginaire sont
les formules de la proposition, $W'' = (1+i)^2a^2W = 2ia^2W$. Les deux équations
se combinent en $A_z(u'' + iv'') = -fv + ifu = if(u + iv)$, c'est-à-dire $A_zW'' = ifW$,
ce qui est vrai puisque $2ia^2A_z = if$ pour $a^2 = f/(2A_z)$. (b)~$W'(0) =
(1+i)aV_0e^{i\pi/4} = \sqrt2\,aV_0\,e^{i\pi/2} = i\sqrt2\,aV_0$~: la tension de surface
$\rho A_zW'(0)$ pointe vers le nord, avec la valeur $\rho A_z\sqrt2\,aV_0 = \rho A_z\sqrt2
\sqrt{f/(2A_z)}\;T/(\rho\sqrt{fA_z}) = T$. (c)~$\int_{-\infty}^0 W\,\mathrm dz =
V_0e^{i\pi/4}/((1+i)a) = V_0/(\sqrt2\,a)$, un nombre réel~: le transport n'a aucune
partie vers le nord, et sa partie vers l'est est $\rho V_0/(\sqrt2\,a) =
T/\sqrt{fA_z}\cdot\sqrt{A_z/f} = T/f$. C'est le théorème~: valeur $T/f$, \ang{90} à
droite d'un vent soufflant vers le nord.
\end{solution}

\begin{solution}{exo:b1:ocean-circulation:11}
(a)~Un pas vers le nord $\mathrm dy = a_\oplus\,\mathrm d\varphi$ change $f = 2\Omega
\sin\varphi$ de $2\Omega\cos\varphi\,\mathrm d\varphi$, donc $\beta = 2\Omega
\cos\varphi/a_\oplus = 2 \times \num{7.292e-5} \times \cos\ang{30}/(\num{6.371e6}) =
\qty{1.98e-11}{m^{-1}.s^{-1}}$. (b)~$R/\beta = \num{e-6}/(\num{1.98e-11}) =
\qty{5.04e4}{m}$, environ \qty{50}{km}. (c)~Intérieur~: $\num{30e6}/(\num{5e6} \times
1000) = \qty{6.0e-3}{m/s}$, \qty{6}{mm/s}~; courant de bord~: $\num{30e6}/(\num{5.04e4}
\times 1000) = \qty{0.60}{m/s}$, cent fois plus rapide, le rapport des largeurs. Le
Gulf Stream, à \qty{1.8}{m/s}, est plus rapide encore~: le courant réel est plus étroit que
cette estimation, son flux est concentré près de la surface, et sa recirculation s'ajoute
à son transport.
\end{solution}

\begin{solution}{exo:b1:ocean-circulation:12}
(a)~Partout où la mer est plus profonde que \qty{1000}{m}, la colonne au-dessus de \qty{1000}{m}
contient au plus \qty{1}{km} d'eau, et ailleurs moins~: le volume au-dessus de
\qty{1000}{m} est au plus $\num{361.84e6} \times 1 = \qty{3.618e8}{km^3}$, donc le
volume en dessous est au moins $\num{1.3324e9} - \num{3.618e8} = \qty{9.71e8}{km^3}$.
(b)~$\num{26e6} \times \num{3.156e7} = \qty{8.21e14}{m^3}$ par an,
\qty{8.21e5}{km^3}~; $\num{9.71e8}/(\num{8.21e5}) \approx 1200$ ans, le même ordre que
les mille ans du tapis roulant. (c)~$(1/2)^{25/12.3} = 0.244$, environ un quart~;
$\num{e-3} \times 25 \times \num{3.156e7} = \qty{7.9e5}{m}$, environ \qty{790}{km}.
Après 25 ans le tritium ne s'est pas étalé plus loin qu'une lente dérive ne le porterait,
alors qu'un quart de lui reste pour être détecté.
\end{solution}

\begin{solution}{pb:b1:ocean-circulation:1}
\begin{enumerate}
  \item Un vent du nord souffle vers le sud, le long de la côte.
  \item $1000\,C_D = 0.60 + 0.071 \times 8 = 1.168$~: $C_D = \num{1.168e-3}$.
  \item $T = 1.3 \times \num{1.168e-3} \times 64 = \qty{0.0972}{Pa}$.
  \item $C_D = \num{1.452e-3}$ et $T = \qty{0.272}{Pa}$~: la tension croît de $2.80$,
        plus que $(12/8)^2 = 2.25$, parce que $C_D$ lui-même augmente avec le vent.
  \item $\Omega = \qty{7.292e-5}{rad.s^{-1}}$~; $f = 2\Omega\sin\ang{38} =
        \qty{8.979e-5}{s^{-1}}$.
  \item $2\pi/f = \qty{7.00e4}{s}$, \qty{19.4}{h}. La force de Coriolis fait tourner un
        courant d'un tour complet en ce temps~; avant que le vent ait soufflé
        environ aussi longtemps, l'eau oscille au lieu de se fixer dans l'équilibre
        permanent du théorème.
  \item \ang{90} à droite d'un vent soufflant vers le sud~: vers l'ouest, le
        large.
  \item $M = T/f = 0.0972/(\num{8.979e-5}) = \qty{1082}{kg.m^{-1}.s^{-1}}$.
  \item $1082/1027 = \qty{1.054}{m^2/s}$ par mètre de côte.
  \item $D_E = 7.6 \times 8/\sqrt{\sin\ang{38}} = \qty{77.5}{m}$.
  \item $1.054/77.5 = \qty{0.0136}{m/s}$, \qty{1.36}{cm/s}.
  \item $V_0 = 0.0127 \times 8/\sqrt{\sin\ang{38}} = \qty{0.129}{m/s}$, \ang{45} à
        droite du vent~: vers le sud-ouest. Les dériveurs donnent environ la moitié
        de la valeur d'Ekman, quelque \qty{0.07}{m/s}.
  \item L'eau qui quitte la côte dans la \omterm{def:b1:ocean-circulation:ekman}{couche d'Ekman} ne peut être remplacée à travers la
        côte, ni par le haut~: elle est remplacée par de l'eau montant d'en dessous de la couche
        près de la côte, comme dans la coupe dessinée dans la
        \cref{sec:b1:ocean-circulation:upwelling}~; les isothermes se courbent vers le haut vers le
        rivage.
  \item $1.054 \times 1000 \times 86400 = \qty{9.11e7}{m^3}$ par jour et par kilomètre de
        côte.
  \item $w = 1.054/(\num{2e4}) = \qty{5.27e-5}{m/s}$, \qty{4.55}{m} par jour.
  \item $77.5/4.55 = 17$ jours.
  \item L'air est refroidi par en dessous par l'eau froide~: la couche près de la surface
        devient froide et stable, l'humidité peut se condenser en brouillard et en stratus bas,
        et sans convection il y a peu de pluie sur la côte.
  \item L'eau remontée est riche en nutriments~; ils nourrissent le phytoplancton, qui
        nourrit le zooplancton et les poissons.
  \item Un vent soufflant vers le nord entraîne un transport \ang{90} à sa droite, vers
        l'est~: l'eau de surface s'accumule contre la côte et s'enfonce (\omterm{def:b1:ocean-circulation:upwelling}{plongée d'eaux})~; la
        remontée s'arrête et de l'eau de surface plus chaude revient à la côte.
  \item $f(\ang{20}) = \qty{4.988e-5}{s^{-1}}$~: $M = \qty{1948}{kg.m^{-1}.s^{-1}}$,
        $\sin\ang{38}/\sin\ang{20} = 1.80$ fois plus.
  \item Dans l'hémisphère Sud le transport est \ang{90} à gauche du
        vent. Pour le pousser vers l'ouest, vers le large, le vent doit souffler vers le nord, venant du
        sud~: dans les deux cas le vent souffle le long de la côte vers l'équateur.
  \item $1000\,C_D = 0.29 + 0.62 + 0.308 = 1.218$~; $T = \qty{0.0396}{Pa}$~; $M =
        \qty{441}{kg.m^{-1}.s^{-1}}$~; le rapport est $0.407$ pour les deux.
  \item Par mètre de côte, $M c_p\,\Delta T = 1082 \times 3992 \times 6 =
        \qty{2.59e7}{W}$, \qty{25.9}{MW}~; le long de \qty{1000}{km}, \qty{2.59e13}{W},
        \qty{25.9}{TW}.
  \item $\num{31.7e6}/1.054 = \qty{3.01e7}{m}$, environ \qty{30000}{km}~: la \omterm{def:b1:ocean-circulation:upwelling}{remontée
        côtière} déplace bien moins d'eau qu'un \omterm{def:b1:ocean-circulation:wbc}{courant de bord ouest}.
  \item $Q = 1.054 \times \num{e6} = \qty{1.05e6}{m^3/s}$~: le transport de remontée
        le long de \qty{1000}{km} de côte est de \qty{1.05}{\sverdrup}.
\end{enumerate}
\end{solution}
